% Generated draft: requires mathematical and source review.
\documentclass[11pt]{article}
\usepackage[margin=25mm]{geometry}
\usepackage{fontspec}
\setmainfont{Latin Modern Roman}
\usepackage{amsmath,amssymb,mathtools,unicode-math}
\setmathfont{Latin Modern Math}
\usepackage{enumitem,needspace,xurl,hyperref}
\setlength{\parindent}{0pt}
\setlength{\parskip}{5pt}
\setlength{\emergencystretch}{3em}
\newcommand{\N}{\mathbb{N}}
\newcommand{\Z}{\mathbb{Z}}
\newcommand{\Q}{\mathbb{Q}}
\newcommand{\R}{\mathbb{R}}
\newcommand{\C}{\mathbb{C}}
\newcommand{\F}{\mathbb{F}}
\newcommand{\A}{\mathbb{A}}
\newcommand{\PP}{\mathbb{P}}
\newcommand{\E}{\mathbb{E}}
\newcommand{\eps}{\varepsilon}
\newcommand{\dd}{\,\mathrm{d}}
\begin{document}
\section*{UnsolvedMath ID: 20003191}
\textbf{Source problem code:} catalogue entry
\subsection{Definitions and mathematical statement}
Let \(S\) be a closed connected oriented surface of genus at least two, \(\Gamma=\pi_1(S)\), and \(\rho_0:\Gamma\to PSL(2,\R)\) a Fuchsian representation. Choose a lift \(\widetilde\rho_0:\Gamma\to SL(2,\R)\). The question of Anna Wienhard asks for embeddings \(\phi:SL(2,\R)\hookrightarrow G\) for which the connected component containing \([\phi\circ\widetilde\rho_0]\) consists entirely of discrete faithful representations. Interpret \(\operatorname{Rep}(\Gamma,G)\) here as the reductive real character space with its usual topology, as in the supplied aid; the original notation does not specify reductive versus raw conjugacy quotient conventions. The damaged diagram simply depicts this composite. Its following page number is omitted.

The source's Comment 5.3 records the split-real irreducible/principal construction as a Hitchin component, and tight embeddings into Hermitian groups as yielding a component of maximal representations. Tightness and maximality involve the appropriate bounded K\"ahler class and Toledo normalization; full target-group hypotheses and theorem references are not given. These known families motivate the broader question without supplying a classification for arbitrary \(G\).

The aid gives a precise parity-safe split example. For \(n\ge2\), let
\[
\phi_n=\operatorname{Sym}^{n-1}:SL(2,\R)\longrightarrow SL(n,\R).
\]
Since \(\phi_n(-I)=(-1)^{n-1}I\), its kernel is trivial for even \(n\) and is \(\{\pm I\}\) for odd \(n\). Thus for every even \(n\), \(\phi_n\) is a literal Lie-group embedding and the component of the reductive \(SL(n,\R)\) character space containing its lifted Fuchsian composite has only discrete faithful representations, according to the aid.

The reported argument projects the component by the finite central map \(SL(n,\R)\to PSL(n,\R)\). Its connected image lies in the projective Hitchin component, whose representations are discrete faithful in the applicable theorem. For a finite central isogeny \(p:\widetilde G\to G\), if \(p\circ\rho\) is discrete faithful, then \(\rho\) is faithful and has discrete image: faithfulness follows from that of the composite, and the preimage of the discrete image under a finite covering is discrete. Hence such a good projective component transfers this property to every connected lifted component mapping into it.

For odd \(n\), the group homomorphism \(\phi_n\) is not an embedding of literal \(SL(2,\R)\), but its surface composite is still faithful: the Fuchsian lift meets \(\{\pm I\}\) only in the identity, because its projection is faithful. Thus it is essential not to confuse the kernel of the ambient homomorphism with the kernel of the surface representation. This supplies explicit examples for the stated reductive-component interpretation and a transfer principle, not an all-target classification or an unqualified identification of \(SL_2\) with \(PSL_2\).
\subsection{Short English statement}
Seek components made entirely of discrete faithful surface-group representations arising from an \(SL_2\) embedding. The supplied split examples are literal embeddings in even dimension; odd dimension has a central kernel even though the surface composite stays faithful. Finite-central lifting preserves the good-component property.
\subsection{Sources}
\begin{itemize}
\item Anna Wienhard, Question 5.2 and Comment 5.3, AIM workshop \emph{Representations of surface groups}, \url{https://aimath.org/WWN/surfacegroups/surfacegroups.pdf}. The comment names the Hitchin and tight Hermitian maximal-representation constructions, without separate bibliography.
\item Local UnsolvedMath registry, permanent ID 20003191, exact code \texttt{catalogue entry}; canonical location \texttt{aim-topology-notes.json}, notes[278], dated 14 August 2026; statement labeled exact. The supplied verification reconstructs the damaged diagram and states the closed oriented genus-at-least-two surface assumptions from the official four-page PDF.
\item Research attempt 1: partial result, problem partially solved at the run, rigor label proved, novelty confidence low; registry status partially solved. The finite-central and parity account is a machine-generated research aid, not peer reviewed; independent verification is required before citation. Full Hitchin/maximal-representation theorem references are not supplied.
\end{itemize}
\end{document}
